\section{Conclusiones}

Las conclusiones deben sintetizar los hallazgos, reconocer las limitaciones de la evidencia y proponer un plan de mejora priorizado. No repita el procedimiento de cada caso: interprete qué demuestra la evidencia, qué no demuestra y qué control reduce el riesgo.

\begin{table}[H]
	\caption{Síntesis de hallazgos y controles por caso}
	\label{tab:sintesis-conclusiones}
	\centering
	\footnotesize
	\setlength{\tabcolsep}{3pt}
	\renewcommand{\arraystretch}{1.2}
	\begin{tabularx}{\textwidth}{@{}c >{\raggedright\arraybackslash}p{0.20\textwidth} Y >{\raggedright\arraybackslash}p{0.22\textwidth} >{\raggedright\arraybackslash}p{0.16\textwidth}@{}}
		\toprule
		\textbf{Caso} & \textbf{Hallazgo principal} & \textbf{Evidencia que lo sustenta} & \textbf{Control recomendado} & \textbf{Limitación} \\
		\midrule
		1 & & & & \\
		2 & & & & \\
		3 & & & & \\
		4 & & & & \\
		5 & & & & \\
		6 & & & & \\
		7 & & & & \\
		8 & & & & \\
		9 & & & & \\
		10 & & & & \\
		\bottomrule
	\end{tabularx}
\end{table}

\subsection{Conclusiones de la Parte I}

\textit{[Redacte entre tres y cinco párrafos. Indique qué patrones se repitieron entre casos, qué controles tuvieron mayor impacto y qué indicadores resultaron más fiables. Diferencie los hallazgos confirmados de las hipótesis no resueltas.]}

\subsection{Conclusiones de la Parte II}

\textit{[Explique qué solución se propuso, por qué es viable, qué requisitos de seguridad y privacidad se priorizaron y qué límites tiene el prototipo. Diferencie el rastreo con consentimiento del rastreo encubierto y justifique por qué este último queda fuera del alcance.]}

\subsection{Limitaciones y trabajo futuro}

\enlargethispage*{2\baselineskip}
\begin{itemize}[leftmargin=*, topsep=2pt, itemsep=2pt]
	\item \textbf{Limitaciones de la evidencia:} tamaño de las capturas, tráfico sintético, ausencia de algunos registros, relojes no sincronizados o datos anonimizados que impiden reproducir el contexto.
	\item \textbf{Limitaciones del prototipo:} cobertura de dispositivos, precisión del GPS, disponibilidad de APIs, consumo de batería y dependencia de servicios de terceros.
	\item \looseness=-1 \textbf{Trabajo futuro:} ampliar la línea base, automatizar la correlación con reglas explicables, evaluar la deriva del modelo, incorporar pruebas de carga y mejorar los mecanismos de consentimiento y borrado.
\end{itemize}

Los Casos 4 y 9 muestran que los prototipos son útiles para priorizar y explicar señales, pero no bastan como mecanismo autónomo de bloqueo. En el Caso 4 se debe ampliar la línea base con más cambios legítimos de carga, vigilar la deriva y proteger los datos de entrenamiento; en el Caso 9 se debe incorporar la telemetría ausente, validar las reglas de correlación en Wazuh y mantener la revisión humana. En ambos casos, la mejora prioritaria es ampliar la evidencia y validar el comportamiento fuera de las capturas controladas.
